\section{Introduction}\label{sec:introduction}

\subsection{Background and Research Gap}

Dementia represents a growing global public health challenge associated with aging populations and increasing pressure on healthcare systems. With over 55 million people affected worldwide and nearly 10 million new cases annually, dementia is among the leading causes of disability and dependency in older adults \citep{who_dementia}. Beyond the substantial emotional, social, and economic burden on patients and caregivers, early recognition and appropriate clinical management remain critical for improving outcomes and care planning.

Primary care clinicians play a central role in dementia care, serving as the first point of contact for patients and families. However, dementia consultations are particularly challenging due to diagnostic uncertainty, the need to integrate multiple assessment dimensions, and significant communication demands. Discussions regarding cognitive decline, loss of autonomy, or care planning may generate fear, denial, or emotional distress among patients and caregivers. Clinicians must balance clinical reasoning with empathetic communication while adapting their approach to sensitive interpersonal dynamics.

While LLMs have enabled various clinical support applications, existing AI systems primarily focus on documentation automation or generalized decision support. A critical gap remains: most tools do not actively guide clinicians through structured consultations or address relational communication challenges specific to dementia care. Concerns regarding hallucinations, explainability, and integration with clinical workflows further limit adoption in sensitive settings.

\subsection{Research Aim and Objectives}\label{sec:aim}

This project addresses the identified gap by developing a guided consultation assistant that combines LLMs with evidence-based communication frameworks to support clinicians throughout dementia consultations.

To achieve this aim, the following objectives were pursued:

\begin{enumerate}
    \item To review existing approaches for AI-assisted clinical support and dementia communication tools, identifying gaps in guidance, evidence grounding, and evaluation (Section~\ref{sec:relatedwork}).
    
    \item To design and implement a retrieval-augmented generation system integrating the Ariadne Labs Essential Communications Toolkit with LLMs (Section~\ref{sec:methods}).
    
    \item To develop structured workflow navigation aligned with the Navigation Map framework across Recognition, Evaluation, and Diagnosis phases.
    
    \item To implement the Stuck Points Framework for identifying and addressing relational communication obstacles.
\end{enumerate}

\subsection{Contributions}

This project makes the following contributions:

\begin{itemize}
    \item A standalone web-based clinical support system providing real-time conversational guidance grounded in evidence-based resources.
    
    \item Integration of the Navigation Map structure enabling structured workflow navigation across dementia consultation phases.
    
    \item A specialized Stuck Points Framework module for relational communication support during sensitive discussions.
\end{itemize}

These contributions demonstrate the potential for combining LLMs with structured communication frameworks to support communication-centered clinical workflows in dementia care.
